\documentclass[12pt]{article}
\usepackage{amsmath,amssymb,amsthm}
\usepackage[margin=1in]{geometry}

\title{Project Euler Problem 498: Remainder of Polynomial Division}
\author{}
\date{}

\begin{document}
\maketitle

\section{Problem Statement}
For positive integers $n$ and $m$, define $F_n(x) = x^n$ and $G_m(x) = (x-1)^m$. Let $R_{n,m}(x)$ be the remainder when dividing $F_n(x)$ by $G_m(x)$, and $C(n,m,d)$ the absolute value of the coefficient of the $d$-th degree term.

Find $C(10^{13}, 10^{12}, 10^4) \bmod 999999937$.

\section{Mathematical Analysis}

\subsection{Change of Variables}
Let $t = x - 1$:
\[
R_{n,m}(x) = \sum_{i=0}^{m-1} \binom{n}{i}(x-1)^i
\]

\subsection{Coefficient Extraction}
\[
[x^d]\, R_{n,m}(x) = \sum_{i=d}^{m-1} \binom{n}{i}\binom{i}{d}(-1)^{i-d}
\]

\subsection{Simplification via Identity}
Using $\binom{n}{i}\binom{i}{d} = \binom{n}{d}\binom{n-d}{i-d}$:
\[
C(n,m,d) = \left|\binom{n}{d} \sum_{k=0}^{m-d-1}\binom{n-d}{k}(-1)^k\right|
\]

\subsection{Lucas' Theorem}
Since $p = 999999937$ is prime and $n > p$, we apply Lucas' theorem to compute binomial coefficients modulo $p$ via base-$p$ digit decomposition.

\section{Complexity}
\begin{itemize}
    \item \textbf{Time:} $O(p)$ for factorial precomputation.
    \item \textbf{Space:} $O(p)$.
\end{itemize}

\section{Answer}

\[
\boxed{472294837}
\]

\end{document}
